%
% ---------------------------------------------------------------
% Copyright (C) 2012-2018 Gang Li
% ---------------------------------------------------------------
%
% This work is the default powerdot-tuliplab style test file and may be
% distributed and/or modified under the conditions of the LaTeX Project Public
% License, either version 1.3 of this license or (at your option) any later
% version. The latest version of this license is in
% http://www.latex-project.org/lppl.txt and version 1.3 or later is part of all
% distributions of LaTeX version 2003/12/01 or later.
%
% This work has the LPPL maintenance status "maintained".
%
% This Current Maintainer of this work is Gang Li.
%
%

\documentclass[
 size=14pt,
 paper=smartboard,  %a4paper, smartboard, screen
 mode=present, 		%present, handout, print
 display=slides, 	% slidesnotes, notes, slides
 style=tuliplab,  	% TULIP Lab style
 pauseslide,
 fleqn,leqno]{powerdot}


%我自己增加的两端对其用
\usepackage{ragged2e}
\renewcommand{\raggedright}{\leftskip=0pt \rightskip=0pt plus 0cm}


\usepackage{cancel}
\usepackage{caption}
\usepackage{stackengine}
\usepackage{smartdiagram}
\usepackage{attrib}
\usepackage{amssymb}
\usepackage{amsmath} 
\usepackage{amsthm} 
\usepackage{mathtools}
\usepackage{rotating}
\usepackage{graphicx}
\usepackage{boxedminipage}
\usepackage{rotate}
\usepackage{calc}
\usepackage[absolute]{textpos}
\usepackage{psfrag,overpic}
\usepackage{fouriernc}
\usepackage{pstricks,pst-3d,pst-grad,pstricks-add,pst-text,pst-node,pst-tree}
\usepackage{moreverb,epsfig,subfigure}
\usepackage{color}
\usepackage{booktabs}
\usepackage{etex}
\usepackage{breqn}
\usepackage{multirow}
\usepackage{natbib}
\usepackage{bibentry}
\usepackage{gitinfo2}
\usepackage{siunitx}
\usepackage{nicefrac}
%\usepackage{geometry}
%\geometry{verbose,letterpaper}
\usepackage{media9}
\usepackage{animate}
%\usepackage{movie15}
\usepackage{auto-pst-pdf}

%\usepackage{breakurl}
\usepackage{fontawesome}
\usepackage{xcolor}
\usepackage{multicol}



\usepackage{verbatim}
\usepackage[utf8]{inputenc}
\usepackage{dtk-logos}
\usepackage{tikz}
\usepackage{adigraph}
%\usepackage{tkz-graph}
\usepackage{hyperref}
%\usepackage{ulem}
\usepackage{pgfplots}
\usepackage{verbatim}
\usepackage{fontawesome}


\usepackage{todonotes}
% \usepackage{pst-rel-points}
\usepackage{animate}
\usepackage{fontawesome}

\usepackage{listings}
\lstset{frameround=fttt,
frame=trBL,
stringstyle=\ttfamily,
backgroundcolor=\color{yellow!20},
basicstyle=\footnotesize\ttfamily}
\lstnewenvironment{code}{
\lstset{frame=single,escapeinside=`',
backgroundcolor=\color{yellow!20},
basicstyle=\footnotesize\ttfamily}
}{}


\usepackage{hyperref}
\hypersetup{ % TODO: PDF meta Data
  pdftitle={Kaggle sildes},
  pdfauthor={Wang Mingxi},
  pdfpagemode={FullScreen},
  pdfborder={0 0 0}
}


% \usepackage{auto-pst-pdf}
% package to show source code

\definecolor{LightGray}{rgb}{0.9,0.9,0.9}
\newlength{\pixel}\setlength\pixel{0.000714285714\slidewidth}
\setlength{\TPHorizModule}{\slidewidth}
\setlength{\TPVertModule}{\slideheight}
\newcommand\highlight[1]{\fbox{#1}}
\newcommand\icite[1]{{\footnotesize [#1]}}

\newcommand\twotonebox[2]{\fcolorbox{pdcolor2}{pdcolor2}
{#1\vphantom{#2}}\fcolorbox{pdcolor2}{white}{#2\vphantom{#1}}}
\newcommand\twotoneboxo[2]{\fcolorbox{pdcolor2}{pdcolor2}
{#1}\fcolorbox{pdcolor2}{white}{#2}}
\newcommand\vpspace[1]{\vphantom{\vspace{#1}}}
\newcommand\hpspace[1]{\hphantom{\hspace{#1}}}
\newcommand\COMMENT[1]{}

\newcommand\placepos[3]{\hbox to\z@{\kern#1
        \raisebox{-#2}[\z@][\z@]{#3}\hss}\ignorespaces}

\renewcommand{\baselinestretch}{1.2}


\newcommand{\draftnote}[3]{
	\todo[author=#2,color=#1!30,size=\footnotesize]{\textsf{#3}}	}
% TODO: add yourself here:
%
\newcommand{\gangli}[1]{\draftnote{blue}{GLi:}{#1}}
\newcommand{\shaoni}[1]{\draftnote{green}{sn:}{#1}}
\newcommand{\gliMarker}
	{\todo[author=GLi,size=\tiny,inline,color=blue!40]
	{Gang Li has worked up to here.}}
\newcommand{\snMarker}
	{\todo[author=Sn,size=\tiny,inline,color=green!40]
	{Shaoni has worked up to here.}}

%%%%%%%%%%%%%%%%%%%%%%%%%%%%%%%%%%%%%%%%%%%%%%%%%%%%%%%%%%%%%%%%%%%%%%%%
% title
% TODO: Customize to your Own Title, Name, Address
%
\title{The Final Report}
\author{
Wang Mingxi
\\
\\Jilin University
\\College of Computer Science and Technology
}
%\date{\gitCommitterDate}
\date{\today} %暂时手写改动

% Customize the setting of slides
\pdsetup{
% TODO: Customize the left footer, and right footer
rf=\href{http://www.tulip.org.au}{
%Last Changed by: \textsc{\gitCommitterName}\ \gitVtagn-\gitAbbrevHash\ (\gitAuthorDate)
Last Changed by: \textsc{Mingxi Wang}\ (\today)
},
cf={flip01-kaggle},
}


\begin{document}

\maketitle

%\begin{slide}{Overview}
%\tableofcontents[content=sections]
%\end{slide}


%%==========================================================================================
%%
\begin{slide}[toc=,bm=]{Overview}
\tableofcontents[content=currentsection,type=1]
\end{slide}
%%
%%==========================================================================================


\section{Research motivation and context}

%%==========================================================================================
%%
\begin{slide}{Project Objectives}

\begin{itemize}
\item Determine whether the disaster reflected in the tweets is real
\end{itemize}

\begin{center}
	\begin{figure}[htbp]
		\includegraphics[scale=0.5]{./pic/1.eps}
	\end{figure}
\end{center}

\end{slide}
%%
%%==========================================================================================

%%==========================================================================================
%%
\begin{slide}{Project background I}
A tweet that says there's a fire in ABA, ablaze means fire, it's the key word in the whole tweet, it means it's a disaster tweet, so it's marked 1. It could also be a key word in "getting people angry," which is disaster\_neutral Twitter, and it should be flagged zero. \\
\begin{center}
	\begin{figure}[htbp]
		\includegraphics[scale=0.8]{./pic/2.eps}
	\end{figure}
\end{center}
\end{slide}
\begin{slide}{Project background II}
	So our task was to build a model that could pick out tweets that were actually disasters and mark them as 1, and distinguish between tweets that were not disasters and tweets that were not disasters and mark them as 0. The test.csv file has 3,263 such tweets waiting to be tagged.
\end{slide}

\section{Research contents and methods}

%%
%%==========================================================================================
\begin{slide}{Import the file AND Discard the same data}
%当有多个Kaggle Subject而只想在目录出现一次时，在后面几次Kaggle Subject前加上[toc=,bm=]
train=pd.read\_csv('./nlp\_getting\_started/train.csv')\\
test=pd.read\_csv('./nlp\_getting\_started/test.csv')\\
\\
df = data\\
df = data.drop\_duplicates().\\
\end{slide}

\begin{slide}{Solve the category imbalance}
\begin{itemize}
	\item  The imbalance of categories (i.e., labels, 0,1) causes the classifier to bias the test set prediction. Too many positive examples in the training set will lead to the model's tendency to predict the test set as positive examples and vice versa. \\
\item Target = 0 has 4,322 tweets and target = 1 has 3,239 tweets. It's about 4 3. It's OK, don't change it.\\
\end{itemize}
\end{slide}

\begin{slide}{Data cleaning}
	def clean\_text(text):\\
	temp = text.lower()          \\                       
	temp = re.sub('n', ' ' , temp)   \\                  
	temp = re.sub('\'', '', temp)     \\                  
	temp = re.sub('', ' ', temp)    \\                   
	temp = re.sub(r'(http|https|pic.)S', ' ', temp)    \\
	temp = re.sub(r'[s]', ' ', temp)       \\         
	return temp
\end{slide}
\begin{slide}{Delete stop word}
def remove\_stopwords(text):\\
	temp = [text for text in text.split() if len(text) > 3]\\
	tokenized\_words = word\_tokenize(text)\\
	temp = [word for word in tokenized\_words if word not in stop\_words]\\
	temp = ' '.join(temp)\\
	return temp\\
\end{slide}
\begin{slide}{Data cleaning}
\begin{itemize}
	\item There is an artificial addition to the dataset of "clean" : the Twitter text after it has been cleaned.\\
\end{itemize}
train['clean'] = train['text'].apply(clean\_text)\\
train['clean'] = train['clean'].apply(remove\_stopwords)\\
train['clean'] = train['text'].apply(clean\_text)\\
train['clean'] = train['clean'].apply(remove\_stopwords)\\
\end{slide}
\begin{slide}{Add artificial features}
			\begin{itemize}
			\item In order to analyze the needs, we should first establish a new combination attribute: that is, the cleaned feature and the keyword are combined with blank space as the new feature.\\
			\end{itemize}
		def combine\_attributes(text, keyword):\\
		var\_list = [text, keyword]\\
		combined = ' '.join(x for x in var\_list if x)\\
		return combined\\
		train.fillna('', inplace = True)\\
		train['combine'] = train.apply(lambda x:\\ combine\_attributes(x['clean'],x['keyword']), axis = 1)\\
		test.fillna('', inplace = True)\\
		test['combine'] = test.apply(lambda x:\\ combine\_attributes(x['clean'],x['keyword']), axis = 1)\\
\end{slide}
\begin{slide}{Data set segmentation}
		\begin{itemize}
			\item The training data were divided according to the ratio of 8 2, with 0.8 of the data training model and 0.2 of the data testing the training effect of the model.\\
		\end{itemize}
		X = train['combine']\\
		y = train['target']\\
		X\_train, X\_test, y\_train, y\_test = train\_test\_split(X, y, train\_size=0.8)\\
\end{slide}

\begin{slide}[toc=,bm=]{Word frequency inverse text frequency (TF IDF) processing}
	\begin{itemize}
	\item The word frequency\_inverse text frequency (TF\_IDF) processing is used to add weight to words in the text.
	\end{itemize}
	from sklearn.feature\_extraction.text import TfidfVectorizer\\
	\\
	vectorizer = TfidfVectorizer()\\
	\\
	X\_train\_vect = vectorizer.fit\_transform(X\_train)\\
	X\_train\_vect\_all = vectorizer.transform(train['clean'])\\
	X\_test\_vect = vectorizer.transform(X\_test)\\
	X\_test\_vect\_all = vectorizer.transform(test['clean'])\\

\end{slide}

\begin{slide}{Visualization heat map I}
import seaborn as sns 
from sklearn.metrics import accuracy\_score, confusion\_matrix 
import matplotlib.pyplot as plt 
conf\_mat = confusion\_matrix(y\_test, y\_pred) 
fig, ax = plt.subplots(figsize=(10,8)) 
sns.heatmap(conf\_mat,cmap="Oranges")  
plt.ylabel('actually',fontsize=18) 
plt.xlabel('predict',fontsize=18) 
\begin{center}
	\begin{figure}[htbp]
		\includegraphics[scale=0.8]{./pic/3.eps}
	\end{figure}
\end{center}
	
\end{slide}
\begin{slide}{Visualization heat map II}
As you can see, there are fewer categories of 1 (target = 1) and more categories of 0. And the prediction was wrong more often when category 0 was classified as category 1 than when 1 was classified as category 0.
\end{slide}


\begin{slide}{The model was created by "support vector machine classification" SVC}
	from sklearn.svm import SVC\\
	from sklearn.svm import LinearSVC\\
	
	clf = SVC(kernel = 'linear')\\
	clf.fit(X\_train\_vect, y\_train)\\
	
	y\_pred = clf.predict(X\_test\_vect)\\
\end{slide}

\section{Conclusion}

\begin{slide}[toc=,bm=]{Evaluation model}
from sklearn.metrics import accuracy\_score\\
accuracy\_score(y\_test, y\_pred)\\
\\
RESULT : 0.8006\\
	
\end{slide}


%%
%%==========================================================================================

%%==========================================================================================
%%

%%
%%==========================================================================================



%%
%%==========================================================================================


%%==========================================================================================
% TODO: Contact Page
\begin{wideslide}[toc=,bm=]{Contact Information}
\centering
\vspace{\stretch{1}}
\twocolumn[
lcolwidth=0.35\linewidth,
rcolwidth=0.65\linewidth
]
{
% \centerline{\includegraphics[scale=.2]{tulip-logo.eps}}
}
{
\vspace{\stretch{1}}
Wang Mingxi\\
College of Computer Science and Technology\\
Jilin University, China
\begin{description}
 \item[\textcolor{orange}{\faEnvelope}] \href{mailto:mxwang@tulip.academy}
 {\textsc{\footnotesize{mxwang@tulip.academy}}}

 \item[\textcolor{orange}{\faHome}] \href{http://www.tulip.org.au}
 {\textsc{\footnotesize{Team for Universal Learning and Intelligent Processing}}}
\end{description}
}
\vspace{\stretch{1}}
\end{wideslide}

\end{document}

\endinput
